\section{Low-Energy Electromagnetic Role Identification}

The preceding sections derive a compact-fiber structural bridge coupling
\[
\alpha_F=\frac1{x_{\rm fiber}},
\]
where \(x_{\rm fiber}\) is the positive root of the compact-fiber bridge polynomial. The remaining question is the physical role of this structural coupling. The identification
\[
\alpha_F=\alpha_{\rm low-energy}
\]
is not made by numerical agreement alone. It is made by role: the compact-fiber electromagnetic bridge is the material/radiative interface, and leading atomic binding reads that bridge through a two-vertex source-response process.

\subsection{Electromagnetic bridge as material/radiative interface}

In the compact-fiber framework, material closure is represented by
\[
F=Z_4\times Z_4\times Z_8,
\]
while radiative/vibrational placement is represented by
\[
V_{\rm rad}.
\]

The electromagnetic interaction is identified as the carrier-level bridge between these two structural regimes:
\[
F
\longleftrightarrow
V_{\rm rad}.
\]

The coupling associated with this bridge is
\[
\alpha_F.
\]

It is not an element-level label depending on a particular pair
\[
(f,v)\in F\times V_{\rm rad}.
\]
It is the scalar bridge normalization of the material/radiative interface itself. This is why the derivation solves for
\[
x=\alpha_F^{-1}
\]
as an effective bridge capacity.

\subsection{Two-vertex source-response structure}

A leading atomic binding observable is not a one-way bridge event. It is a closed source-response process:
\[
\text{material source}
\rightarrow
\text{radiative/electromagnetic bridge}
\rightarrow
\text{material response}.
\]

The first bridge vertex carries material structure into the electromagnetic/radiative bridge. The second bridge vertex returns the electromagnetic response to the material sector. Thus the leading atomic source-response scale contains two bridge factors:
\[
\alpha_F\cdot\alpha_F=\alpha_F^2.
\]

Therefore the compact-fiber leading atomic scale has the form
\[
E_{\rm atomic,F}
=
C_{\rm mass/readout}\alpha_F^2,
\]
where \(C_{\rm mass/readout}\) contains dimensional mass, action, and finite-readout matching factors. The finite-readout and matching-correction layer is treated separately in \cite{Wells2026FiniteReadout}. The load-bearing structural fact is the coupling power:
\[
E_{\rm atomic,F}\propto \alpha_F^2.
\]

\subsection{Correspondence with Hartree/Rydberg scaling}

In conventional low-energy atomic physics, the Hartree energy is
\[
E_h=\alpha^2m_ec^2.
\]
The Rydberg constant is
\[
R_\infty=\frac{\alpha^2m_ec}{2h}.
\]

Both expressions show the same structural role:
\[
\text{atomic binding scale}\propto\alpha^2.
\]

The compact-fiber source-response structure gives
\[
E_{\rm atomic,F}
=
C_{\rm mass/readout}\alpha_F^2.
\]

Thus the compact-fiber bridge coupling occupies the same low-energy role as the conventional fine-structure constant:
\[
\alpha_F
\leftrightarrow
\alpha_{\rm low-energy}.
\]

This is the physical role identification used in this paper.

\subsection{Boundary of the role identification}

The role identification claims that \(\alpha_F\) is the compact-fiber structural counterpart of the low-energy electromagnetic coupling.

It does not independently derive the dimensional anchors or correction pipelines that appear in particular empirical routes. Those are route-specific readout problems.

The alpha theorem derived in this paper supplies the structural electromagnetic bridge coupling. Other constants and correction layers determine how this coupling appears in specific measurement pipelines.

\subsection{Relation to the Planck/action bridge}

Some empirical routes to \(\alpha\), especially Rydberg/recoil and
\(h/m\)-based routes, involve the action anchor
\[
h.
\]

The compact-fiber program treats this as a separate structural problem.
The Planck/action and gravitational-coupling analysis is developed in a
separate constants-sector paper \cite{Wells2026PlanckGravity}. That
work gives the Planck/action bridge its own structural content,
schematically of the form
\[
h=\frac{m_n c S_n^2}{R_{\rm eff}},
\]
where \(R_{\rm eff}\) is a dimensionless compact-fiber structural factor
and the remaining quantities are dimensional anchors connecting the
structural/natural-unit system to conventional units.

The Planck/action bridge was already identified as a structural scaling
problem in the radiation-sector development \cite{Wells2026Radiation};
the later Planck/gravity analysis treats that bridge as a separate
theorem path rather than as part of the alpha fixed-point theorem.

This Planck/action result is relevant here only to empirical route
correspondence, because several operational extractions of \(\alpha\)
combine electromagnetic coupling, action, mass, and frequency/length
readouts. The separation of layers is simple: the alpha theorem derives
\(\alpha_F\); the Planck/action theorem supplies separate structural
action content; and empirical route analysis describes how operational
determinations combine electromagnetic coupling, action, mass, charge,
and readout conventions.

This paper uses the Planck/action bridge only as context for empirical
route correspondence. It does not merge the Planck derivation into the
alpha fixed-point theorem.

\subsection{Numerical agreement as consistency check}

The compact-fiber structural prediction is
\[
\alpha_F^{-1}=x_{\rm fiber}.
\]

The numerical agreement with the empirical low-energy fine-structure constant supports the role identification
\[
\alpha_F=\alpha_{\rm low-energy}.
\]

However, the numerical agreement is not the sole basis for the identification. The primary basis is structural: \(\alpha_F\) is the material/radiative bridge coupling, and leading atomic binding reads
\[
\alpha_F^2.
\]

The empirical comparison therefore functions as a consistency check on the role identification, not as the premise from which the identification is made.